% !TEX root = ./index.tex

\maketitle
\tableofcontents

\newpage

\section{Critique of AI-Generated Introduction}

\subsection{Structure}

The introduction begins with a sentence about the Black Summer bushfires, which is followed by a sentence about bushfires in Australia. The first sentence should be relocated such that the direction of focus is from general to specific.

Also prepend a sentence about bushfires beyond Australia to provide additional context.

\subsection{Main Elements}

The introduction includes the broader issue and previous research, but no hypotheses or method. Also, the introduction does not mention any scientific debate, and only describes what was discovered.

The hypotheses and an outline of the method should be appended to the introduction, and any debate on the issue of the report should be stipulated.

\subsection{Use of Scientific Literature}

There are discrepancies in the references of the introduction.

For example, the details in the (Andersen et al., 2024) reference are all wrong, and the (York et al., 2018) reference does not exist at all. There are also minor errors in the statistics taken.

The references must be updated with the correct details and used in the right way.

\subsection{Writing Style}

The writing style of the introduction is sound, being objective and written in the fourth person, however ideas are not logically separated between paragraphs and the present tense is used frequently. It can be improved by adjusting paragraph breaks and using passive past tense where possible.

\subsection{Relevance to Experiment}

Without a hypothesis or method, the introduction is not relevant to the experiment performed. Additionally, the aim to \emph{clarify the ecological consequences of [these] fire regimes \dots} is not the aim of the experiment. A hypothesis and method should be added for context, and the aim must be changed to align with the experiment.

\subsection{Title}

The introduction should include a title that summarises the findings of the report.

\break

\section{Effect of Low- and Moderate-Severity Bushfires on the Biodiversity of Leaf-Litter Invertebrates}

\subsection{Introduction}

Extreme wildfires are occuring more frequently across the world and with greater intensity due to climate change, with burn areas 35 times more in the Chiquitano dry forests and 25 times more in Southern California than estimated without climate change \parencite{kelley-et-al-2025}.
The scale and severity of bushfires in Australia particularly --- driven by prolonged drought and record-breaking temperatures --- exceeded historical norms and raised urgent questions about their ecological consequences \parencite{boer-et-al-2020}.

The 2019--2020 Black Summer bushfires were among the most extensive and ecologically disruptive fire events in Australia’s history, burning more than 24 million hectares and affecting a wide range of ecosystems \parencite{royal-commission-2020}.
During the Black Summer fires, an estimated 60 billion soil and leaf-litter invertebrates were killed in temperate rainforests alone, largely due to the combustion or desiccation of the litter layer that forms their primary habitat \parencite{gibb-et-al-2023}.
While high-severity fires typically resulted in near-complete loss of litter fauna, low- and moderate-severity fires can still reduce litter depth, alter moisture regimes, and disrupt thermal buffering, leading to declines in abundance and shifts in community composition \parencite{brennan-et-al-2009}.
These fires, which typically left portions of the litter layer and canopy intact, are often assumed to be less damaging than high-severity burns.
However, it is suggested in recent studies that even these lower-intensity fires can substantially alter habitat structure and biodiversity, particularly for ground-dwelling invertebrates that depend on leaf-litter for shelter, food, and moisture regulation \parencite{york-2000, andrew-et-al-2000}.

Leaf-litter invertebrates --- including detritivores, micro-predators, and soil engineers --- play essential roles in decomposition, nutrient cycling, and soil formation.
Their sensitivity to microhabitat changes makes them valuable indicators of ecosystem disturbance \parencite{maleque-et-al-2009}, and their role as detrivores can contribute to fire suppression \parencite{gibb-et-al-2023}.

The ecological impacts of these fires on such invertebrates are complex, taxon-specific, and poorly understood \parencite{henry-et-al-2022}.
Some groups, such as ants and certain beetles, may have recolonised rapidly due to high mobility or fire-adapted traits \parencite{andersen-et-al-2012}.
Others --- particularly taxa with limited dispersal ability, narrow microhabitat requirements, or reliance on moist litter microclimates --- showed slower recovery or sustained population declines \parencite{brennan-et-al-2009}.
It is even hypothesised that the impacts of fire on invertebrate populations have been irreversible in some cases \parencite{holz-et-al-2015}.
Understanding how low and moderate fire severity has influenced these communities is therefore critical for conserving forest ecosystems and informing fire-management strategies that balance hazard reduction with biodiversity conservation.

We aimed to determine the effect of low- to moderate-severity fire on the species richness, abundance, and biodiversity of leaf-litter invertebrates.
Samples of leaf-litter were collected from two sites in Conjola National Park which were both affected by the Black Summer bushfires: Site A experienced low-severity fire and Site B experienced moderate-severity fire.
The samples were sorted through by hand and with a microscope.
Invertebrates present in the samples were identified and counted according to lowest recognisable taxonomic unit, then the average species richness, average species abundance, and Simpson's Diversity Index was calculated for each site and compared.

If the species richness and abundance of invertebrate taxa are more affected by higher severity fire, then the species richness and abundance observed in Site A will be greater than that in Site B. Also, if bushfires change the surrounding environment to favour certain taxa, then the biodiversity in Site A will be greater than that in Site B.

\break

\subsection{Methods}

We collected leaf-litter samples on the 28th March 2026 from two sites that both experienced some degree of impact in the NSW ‘Black Summer’ 2019-2020 bushfires, a low-severity burn site (35”14’28°S 150”31’04°E) and a moderate- severity burn site (35”13’53°S 150”29’47°E). We determined fire severity at each site using the Google Earth Engine Burnt Area Map within the SEED map portal (NSW Department of Planning and Environment, 2020). We defined Low Severity as class 1 – minimal change post fire, and Moderate Severity as class 3 – partially affected canopy vegetation. Sites were both part of the Currowan fire in Conjola National Park, and both classified as sclerophyll forest.

We collected leaf-litter samples from low-severity (n=166) and moderate-severity (n=160) sites. Approximately 2L of ground leaf-litter was removed by hand to a level of 0.5 cm below the soil surface, bagged, and returned to the laboratory at the University of Sydney. Between 30st March – 3rd April 2026, we placed each sample into a plastic tray and the leaf-litter sorted macroscopically (by hand) and microscopically using a dissecting microscope. We identified and counted all invertebrates to lowest recognisable taxonomic unit.

We calculated average species richness and abundance and their associated standard errors in Excel. For both average richness, and average abundance, we determined if there were differences between low-severity burn sites and moderate\\-severity burn sites using two one-tailed t-tests in Excel.

We used Simpson's Index of Diversity to measure the species diversity of the low-severity and moderate-severity burn sites. The index for each type of site was calculated with the formula
\[ D = 1 - \frac{1}{N^2} \sum^n n_i^2, \]
where $n_i$ was the abundance of each taxon, and $N$ was the total species abundance.

For the low-severity burn sites, the index was
\[
    D = 1 - \frac{1}{2718^2} \left(168^2 + 416^2 + 67^2 + \ldots + 2^2\right) \approx 0.918,
\]

and for the moderate-severity burn sites, the index was
\[
    D = 1 - \frac{1}{2480^2} \left(329^2 + 363^2 + 55^2 + \ldots + 0^2\right) \approx 0.915.
\]

\break

\subsection{Results}

The average species abundance of sampled invertebrates (\cref{fig:abundance}) was 79.48 for Site A (s.e = 22.97), and 62.05 for Site B (s.e = 19.36). There was no significant difference between the two sites (t = 1.161, df = 324, p = 0.246).

\begin{figure}[htp]
    \centering
    \includegraphics[width=0.85\textwidth]{abundance.jpg}
    \caption{Average species abundance of invertebrates from Site A (low-severity fire) and Site B (moderate-severity fire).}
    \label{fig:abundance}
\end{figure}

The average species richness of sampled invertebrates (\cref{fig:species-richness}) was 6.17 for Site A (s.e = 0.18), and 6.11 for Site B (s.e = 0.21). Again, there was no significant difference between the two sites (t = 0.822, df = 324, p = 0.412).

\begin{figure}[htp]
    \centering
    \includegraphics[width=0.85\textwidth]{species-richness.jpg}
    \caption{Average species richness of invertebrates from Site A (low-severity fire) and Site B (moderate-severity fire).}
    \label{fig:species-richness}
\end{figure}

Simpson's Index of Diversity was calculated for each site (\cref{fig:diversity}): the index was 0.918 for Site A and 0.915 for Site B.

\begin{figure}[htp]
    \centering
    \includegraphics[width=0.85\textwidth]{diversity.jpg}
    \caption{Simpson's Index of Diversity calculated for invertebrates from Site A (low-severity fire) and Site B (moderate-severity fire).}
    \label{fig:diversity}
\end{figure}

\pagebreak

\subsection{Discussion}

\break

\section{References}

\printbibliography[heading=none]
